\documentclass[11pt,a4paper]{article}

\usepackage[utf8]{inputenc}
\usepackage[T1]{fontenc}
\usepackage[english]{babel}
\usepackage[a4paper,margin=2.3cm]{geometry}
\usepackage{enumitem}
\usepackage{xcolor}
\usepackage{titlesec}
\usepackage{parskip}
\usepackage{hyperref}
\usepackage{tabularx}
\usepackage{booktabs}
\usepackage{array}



\titleformat{\section}
  {\color{primaryblue}\Large\bfseries}
  {\thesection}{0.6em}{}
\titleformat{\subsection}
  {\color{darkgray}\large\bfseries}
  {\thesubsection}{0.5em}{}

\setlist[itemize]{leftmargin=1.5em,itemsep=2pt,topsep=4pt}

\title{\vspace{-1cm}\textbf{Projet Traveling}\\[0.3em]
  \large Rapport des fonctionnalités réalisées}
\author{ali .... \and BOou ...}

\date{March 2026}


\begin{document}

\maketitle



\vspace{0.5em}
\noindent
\textbf{Dépôt GitLab :} \url{https://gitlab.etu.umontpellier.fr/e20190009868/travelsharepath}

\vspace{1em}

% ============================================================

\section{Travel Share}
\section{TravelPath}

\subsection*{Stack technique}

\begin{itemize}
  \item \textbf{Frontend} : Android natif en Java
  \item \textbf{Backend} : Node.js / Express / MongoDB
  \item \textbf{Architecture} : MVVM (ViewModel + LiveData), Retrofit2, Navigation Component
\end{itemize}

\subsection{Authentification}

\begin{itemize}
  \item Inscription avec nom complet, email et mot de passe .
  \item Connexion avec un token JWT persisté dans les SharedPreferences.
  \item Déconnexion avec nettoyage complet de la session.
  \item L'écran d'accueil adapte son interface selon l'état de connexion.
\end{itemize}

\subsection{Création d'un parcours}

Le formulaire de création est organisé en deux niveaux.

\textbf{Niveau parcours :}
\begin{itemize}
  \item Titre, description, ville.
  \item Météo : Ensoleillé, Nuageux, Pluvieux, Tout temps.
  \item Niveau d'effort : Facile, Moyen (Actif), Difficile.
\end{itemize}

\textbf{Niveau étape (ajout multiple) :}
\begin{itemize}
  \item Nom du lieu, géocodé automatiquement via Nominatim/OSM.
  \item Type : Culture, Découverte, Restauration, Loisirs.
  \item Durée (minutes), budget (€), description.
  \item Photos multiples ajoutées depuis la galerie.
  \item Réordonnancement par glisser-déposer.
  \item Édition inline de chaque étape.
  \item Compteurs en temps réel (durée totale, budget total, nombre d'étapes).
\end{itemize}

Un nouveau parcours est créé en statut \textit{Brouillon}.

\subsection{Recherche de parcours}

\begin{itemize}
  \item Formulaire de filtres : ville, budget (slider), durée (slider), niveau d'effort,
    types d'activités (multi-sélection).
  \item Seuls les parcours \textit{publiés} apparaissent dans les résultats.
  \item Filtrage côté serveur (ville, budget max, durée max, effort,types d'activités, metéo) 
\end{itemize}

\subsection{Interactions sociales}

\textbf{Aimer un parcours} :
\begin{itemize}
  \item Bouton cœur avec compteur visible sur la fiche détail et dans les listes.
  \item Toggle like/unlike confirmé par le serveur (pas d'optimistic update).
  \item Le compteur se met à jour après confirmation serveur.
\end{itemize}

\textbf{Mettre en favori} :
\begin{itemize}
  \item Icône étoile sur la fiche détail.
  \item Liste « Mes Favoris  » dédiée avec les mêmes actions disponibles
    (like, dé-favori, régénérer).
  \item État favori synchronisé via \texttt{GET /auth/me}.
\end{itemize}

\subsection{Publication d'un parcours}

\begin{itemize}
  \item Bouton « Publier  » visible uniquement sur les parcours en statut Brouillon.
  \item Passage à statut \textit{Publié} côté serveur.
  \item La liste « Mes parcours  » se rafraîchit automatiquement après publication.
\end{itemize}

\subsection{Détail d'un parcours --- vue Itinéraire}

\textbf{En-tête :}
\begin{itemize}
  \item Titre, budget total, durée totale, distance estimée, effort, météo, date.
  \item Jusqu'à 4 photos d'aperçu (1\textsuperscript{re} photo de chaque étape).
\end{itemize}

\textbf{Par étape :}
\begin{itemize}
  \item Numéro et nom du lieu.
  \item Badge de type coloré : Culture (indigo), Restauration (orange),
    Loisirs (vert), Découverte (bleu).
  \item Description, durée, budget.
  \item Niveau d'effort coloré (vert / orange / rouge).
  \item Galerie photos swipeable horizontalement.
  \item Temps de marche vers l'étape suivante (si renseigné).
\end{itemize}

\subsection{Détail d'un parcours --- vue Carte}

\begin{itemize}
  \item Carte Leaflet.js intégrée dans une WebView, avec tuiles CartoDB.
  \item Marqueurs numérotés (cercles rouges) pour chaque étape.
  \item Géocodage à la volée des étapes sans coordonnées GPS.
  \item Tracé d'abord en pointillés bleus, puis remplacé par l'itinéraire routier
    réel calculé via OSRM.
\end{itemize}

\subsection{Autres fonctionnalités}

\begin{tabularx}{\linewidth}{@{}>{\bfseries}lX@{}}
\toprule
Mes Parcours & Liste avec actions : voir, éditer, publier, supprimer (avec confirmation). \\
\midrule
Édition / Régénération & Propriétaire $\rightarrow$ mise à jour ; autre utilisateur
  $\rightarrow$ fork (copie). \\
\midrule
Export PDF & PDF multi-pages stylé : couverture, statistiques, description, une carte
  par étape. \\
\midrule
Upload photos & Multipart vers le backend, stocké dans \texttt{/backend/uploads/}. \\
\midrule
Profil & Nom, email, statistiques (parcours créés / publiés), lien vers création. \\
\bottomrule
\end{tabularx}

\vspace{1em}




% ============================================================
\section{Passerelle Traveling}

\textit{À discuter et à compléter en binôme.}

L'objectif de la passerelle est de faire communiquer les deux applications afin
qu'elles forment un système cohérent. L'architecture envisagée s'appuie sur les
mécanismes de communication inter-applications d'Android :

\begin{itemize}
  \item \textbf{Intents} : navigation entre écrans des deux applications.
  \item \textbf{Content Providers} : exposition de données structurées
    (photos d'un lieu, étapes d'un parcours).
  \item \textbf{Services AIDL} : appels de méthodes distantes pour les calculs
    (par exemple, génération d'un parcours à partir des photos populaires d'un lieu).
\end{itemize}

\subsection*{Scénarios envisagés}

\begin{itemize}
  \item Depuis une photo dans TravelShare $\rightarrow$ générer un parcours TravelPath
    incluant ce lieu.
  \item Depuis une étape d'un parcours TravelPath $\rightarrow$ consulter les photos
    associées au lieu sur TravelShare.
\end{itemize}

\end{document}
